% GENERATED FILE. Edit canonical lessons, then run npm run generate:books.
% canonical-source-hash: fnv1a64:54e801075a0011d2
% canonical-lessons: HI-C436-bud, HI-C436-megh, HI-C436-bhap, HI-C436-renkot, HI-C436-barasna

\chapter{Drops and Showers}
\label{ch:hi-drops-and-showers}
\hlchaptermodality{436}

\begin{chapteropening}
\textbf{By the end of this chapter:} \emph{I can talk about drops and spells of rain, steam, a raincoat and rain pouring down.}

Complete the last lesson of chapter 436: I can talk about drops and spells of rain, steam, a raincoat and rain pouring down.
\end{chapteropening}

\section[\texorpdfstring{\hi{बूँद} (bū̃d)}{बूँद (bū̃d)}]{\hi{बूँद} (bū̃d) — a drop}

\label{lesson:HI-C436-bud}



\begin{quote}
Before the new one: say the Hindi for a handkerchief, then the Hindi for to huddle up.
\end{quote}

\subsection*{You'll want to know: \hi{बूँद}}
\textbf{\hi{बूँद}} — \emph{bū̃d} — \textquotedblleft{}a drop\textquotedblright{}.

\begin{quote}
\hi{छत} \hi{से} \hi{पानी} \hi{की} \hi{बूँदें} \hi{टपक} \hi{रही} \hi{थीं।}
\end{quote}

\begin{quote}
\emph{chat se pānī kī bū̃dẽ ṭapak rahī thīṁ.}
\end{quote}

\begin{quote}
Drops of water were dripping from the roof.
\end{quote}

Feminine.

\begin{cousinweb}
\textbf{\hi{बूँद}} (\emph{bū̃d}) is a tadbhava of Sanskrit \emph{bindu}, \textquotedblleft{}a drop, a spot\textquotedblright{} (CDIAL 9240), through a form \emph{*bundu}; CDIAL lists Hindi \emph{bū̃d}, \textquotedblleft{}drop\textquotedblright{}. The dot on a forehead, \hi{बिंदी}, comes from the same word.
\end{cousinweb}

\begin{grammarlens}[title={What you've built}]
One more word for getting around.
\end{grammarlens}

\subsection*{Your turn}
\begin{itemize}
\raggedright
  \item \emph{Say it:} \emph{bū̃d}
  \item \emph{Say it:} \emph{bū̃d}, once more
  \item \emph{Say it:} \emph{chat se pānī kī bū̃dẽ ṭapak rahī thīṁ.} — Drops of water were dripping from the roof.
  \item \emph{Recall:} say the Hindi for a handkerchief, then the Hindi for to huddle up, then say \emph{bū̃d} again
\end{itemize}

\begin{tcolorbox}[breakable,colback=teal!4,colframe=teal!35!black,title={Before you move on}]
What does \hi{बूँद} mean? (\textquotedblleft{}A drop\textquotedblright{}.) Say it once more.
\end{tcolorbox}
\section[\texorpdfstring{\hi{मेघ} (megh)}{मेघ (megh)}]{\hi{मेघ} (megh) — a cloud (literary)}

\label{lesson:HI-C436-megh}



\begin{quote}
Before the new one: say the Hindi for to huddle up, then the Hindi for a drop.
\end{quote}

\subsection*{You'll want to know: \hi{मेघ}}
\textbf{\hi{मेघ}} — \emph{megh} — \textquotedblleft{}a cloud (literary)\textquotedblright{}.

\begin{quote}
\hi{पुराने} \hi{गानों} \hi{में} \hi{मेघ} \hi{बहुत} \hi{आते} \hi{हैं।}
\end{quote}

\begin{quote}
\emph{purāne gānoṁ meṁ megh bahut āte haiṁ.}
\end{quote}

\begin{quote}
Clouds come up a lot in old songs.
\end{quote}

Masculine.

\begin{cousinweb}
\textbf{\hi{मेघ}} (\emph{megh}) is a tatsama, Sanskrit \emph{megha}, \textquotedblleft{}a cloud\textquotedblright{}, literally \textquotedblleft{}the sprinkler\textquotedblright{} (Monier-Williams); poems and songs of the rains are full of it, while talk says \hi{बादल}.
\end{cousinweb}

\begin{grammarlens}[title={What you've built}]
One more word for getting around.
\end{grammarlens}

\subsection*{Your turn}
\begin{itemize}
\raggedright
  \item \emph{Say it:} \emph{megh}
  \item \emph{Say it:} \emph{megh}, once more
  \item \emph{Say it:} \emph{purāne gānoṁ meṁ megh bahut āte haiṁ.} — Clouds come up a lot in old songs.
  \item \emph{Recall:} say the Hindi for to huddle up, then the Hindi for a drop, then say \emph{megh} again
\end{itemize}

\begin{tcolorbox}[breakable,colback=teal!4,colframe=teal!35!black,title={Before you move on}]
What does \hi{मेघ} mean? (\textquotedblleft{}A cloud (literary)\textquotedblright{}.) Say it once more.
\end{tcolorbox}
\section[\texorpdfstring{\hi{भाप} (bhāp)}{भाप (bhāp)}]{\hi{भाप} (bhāp) — steam; vapour}

\label{lesson:HI-C436-bhap}



\begin{quote}
Before the new one: say the Hindi for a drop, then the Hindi for a cloud (literary).
\end{quote}

\subsection*{You'll want to know: \hi{भाप}}
\textbf{\hi{भाप}} — \emph{bhāp} — \textquotedblleft{}steam; vapour\textquotedblright{}.

\begin{quote}
\hi{चाय} \hi{से} \hi{भाप} \hi{उठ} \hi{रही} \hi{थी।}
\end{quote}

\begin{quote}
\emph{cāy se bhāp uṭh rahī thī.}
\end{quote}

\begin{quote}
Steam was rising from the tea.
\end{quote}

Feminine.

\begin{cousinweb}
\textbf{\hi{भाप}} (\emph{bhāp}) is a tadbhava of Sanskrit \emph{bāṣpa}, \textquotedblleft{}tears\textquotedblright{} and also \textquotedblleft{}vapour\textquotedblright{} (CDIAL 9223), which lists Hindi \emph{bhāp}, \textquotedblleft{}steam\textquotedblright{}.
\end{cousinweb}

\begin{grammarlens}[title={What you've built}]
One more word for getting around.
\end{grammarlens}

\subsection*{Your turn}
\begin{itemize}
\raggedright
  \item \emph{Say it:} \emph{bhāp}
  \item \emph{Say it:} \emph{bhāp}, once more
  \item \emph{Say it:} \emph{cāy se bhāp uṭh rahī thī.} — Steam was rising from the tea.
  \item \emph{Recall:} say the Hindi for a drop, then the Hindi for a cloud (literary), then say \emph{bhāp} again
\end{itemize}

\begin{tcolorbox}[breakable,colback=teal!4,colframe=teal!35!black,title={Before you move on}]
What does \hi{भाप} mean? (\textquotedblleft{}Steam; vapour\textquotedblright{}.) Say it once more.
\end{tcolorbox}
\section[\texorpdfstring{\hi{रेनकोट} (renkoṭ)}{रेनकोट (renkoṭ)}]{\hi{रेनकोट} (renkoṭ) — a raincoat}

\label{lesson:HI-C436-renkot}



\begin{quote}
Before the new one: say the Hindi for a cloud (literary), then the Hindi for steam.
\end{quote}

\subsection*{You'll want to know: \hi{रेनकोट}}
\textbf{\hi{रेनकोट}} — \emph{renkoṭ} — \textquotedblleft{}a raincoat\textquotedblright{}.

\begin{quote}
\hi{बारिश} \hi{में} \hi{मैं} \hi{रेनकोट} \hi{पहनता} \hi{हूँ।}
\end{quote}

\begin{quote}
\emph{bāriś meṁ maiṁ renkoṭ pahantā hūṁ.}
\end{quote}

\begin{quote}
I wear a raincoat in the rain.
\end{quote}

Masculine.

\begin{cousinweb}
\textbf{\hi{रेनकोट}} (\emph{renkoṭ}) is borrowed from English \emph{raincoat}.
\end{cousinweb}

\begin{grammarlens}[title={What you've built}]
One more word for getting around.
\end{grammarlens}

\subsection*{Your turn}
\begin{itemize}
\raggedright
  \item \emph{Say it:} \emph{renkoṭ}
  \item \emph{Say it:} \emph{renkoṭ}, once more
  \item \emph{Say it:} \emph{bāriś meṁ maiṁ renkoṭ pahantā hūṁ.} — I wear a raincoat in the rain.
  \item \emph{Recall:} say the Hindi for a cloud (literary), then the Hindi for steam, then say \emph{renkoṭ} again
\end{itemize}

\begin{tcolorbox}[breakable,colback=teal!4,colframe=teal!35!black,title={Before you move on}]
What does \hi{रेनकोट} mean? (\textquotedblleft{}A raincoat\textquotedblright{}.) Say it once more.
\end{tcolorbox}
\section[\texorpdfstring{\hi{बरसना} (barasnā)}{बरसना (barasnā)}]{\hi{बरसना} (barasnā) — to pour (of rain); to rain down}

\label{lesson:HI-C436-barasna}



\begin{quote}
Before the new one: say the Hindi for steam, then the Hindi for a raincoat.
\end{quote}

\subsection*{You'll want to know: \hi{बरसना}}
\textbf{\hi{बरसना}} — \emph{barasnā} — \textquotedblleft{}to pour (of rain); to rain down\textquotedblright{}.

\begin{quote}
\hi{रात} \hi{भर} \hi{बादल} \hi{बरसे।}
\end{quote}

\begin{quote}
\emph{rāt bhar bādal barse.}
\end{quote}

\begin{quote}
It poured all night long.
\end{quote}

Past (it …): \hi{वह} \hi{बादल} \hi{ज़ोर} \hi{से} \hi{बरसा}, \emph{vah bādal zor se barsā}.

\begin{cousinweb}
\textbf{\hi{बरसना}} (\emph{barasnā}) is a tadbhava of Sanskrit \emph{varṣati}, \textquotedblleft{}it rains\textquotedblright{} (CDIAL 11394), which lists Hindi \emph{barasnā}; \hi{वर्षा}, rain, is the same root.
\end{cousinweb}

\begin{grammarlens}[title={What you've built}]
One more word for getting around.
\end{grammarlens}

\subsection*{Your turn}
\begin{itemize}
\raggedright
  \item \emph{Say it:} \emph{barasnā}
  \item \emph{Say it:} \emph{barasnā}, once more
  \item \emph{Say it:} \emph{rāt bhar bādal barse.} — It poured all night long.
  \item \emph{Recall:} say the Hindi for steam, then the Hindi for a raincoat, then say \emph{barasnā} again
\end{itemize}

\begin{tcolorbox}[breakable,colback=teal!4,colframe=teal!35!black,title={Before you move on}]
What does \hi{बरसना} mean? (\textquotedblleft{}To pour (of rain); to rain down\textquotedblright{}.) Say it once more.
\end{tcolorbox}
